While all configurations maintain high coverage during the initial static phase, they demonstrate divergent behaviors during continuous dynamic tracking. Configuration~B heavily prioritizes density ($w_d=0.7$), leading to excessive physical displacement (117.1~m) as the swarm overfits to temporary crowd clusters and needlessly chases crowd dispersal. Furthermore, relaxing the anti-clustering penalty ($\gamma=0.1$) in Configuration~C increases the likelihood of multi-AP interference due to proximity, degrading dynamic coverage to 91.6\%. Configuration~A provides a moderate movement penalty ($\beta=0.15$), but this leads to unnecessary spatial exploration, incurring a high dynamic displacement (113.2~m) without yielding optimal coverage (91.0\%).

Conversely, \textbf{Configuration~D} applies a high movement penalty ($\beta = 0.25$). While originally hypothesized to restrict spatial exploration, the dynamic simulations reveal that this high penalty successfully forces the swarm to rely on efficient electronic beam steering. Under deep geometric blockages (e.g., transitions at $t_0$ and $t_2$), Configuration~D prevents erratic physical relocation, allowing APs to settle into topologically stable, highly efficient positions. Consequently, it minimizes total dynamic displacement to just 42.2~m while sustaining an excellent 95.2\% average dynamic coverage. We therefore adopt \textbf{Configuration~D} as our default parameters across dynamic multi-step evaluations (Section~V-D) because it provides the optimal balance of high coverage stability and minimal mechanical overhead.
